% Solution: Gibbons (1992) Problem 1.2. Input by \solitem.
% Reused from lista-reais-lectures-01-04-solutions (Gibbons 1.2), deepened. Check: check_A1.py.
\begin{exsol}{Gibbons Problem 1.2}{48}{Foundations}{2}
\asked Which strategies of the $3\times3$ game survive IESDS, and which profiles are pure-strategy
Nash equilibria?

\method IESDS (Gibbons \S1.1.B, pp.~4--7; Lecture~1, slides~28--29) and pure-strategy Nash
equilibrium (Gibbons \S1.1.C, pp.~7--12; Lecture~2, slide~14), found by marking best responses
(Lecture~2, slide~10). Gibbons's Proposition~B (p.~12; Lecture~2, Proposition~1, slide~16): every
Nash equilibrium survives IESDS, so after the elimination only the surviving cells need to be
searched. The full marking below confirms it independently.

\solution
The matrix as printed (the top-left cell reads ``2.0'' in the scan; it is $(2,0)$):
\begin{center}
\begin{tabular}{r|c|c|c|}
\multicolumn{1}{r}{1 $\backslash$ 2} & \multicolumn{1}{c}{$L$} & \multicolumn{1}{c}{$C$} & \multicolumn{1}{c}{$R$}\\ \cline{2-4}
$T$ & $2,\ 0$ & $1,\ 1$ & $4,\ 2$ \\ \cline{2-4}
$M$ & $3,\ 4$ & $1,\ 2$ & $2,\ 3$ \\ \cline{2-4}
$B$ & $1,\ 3$ & $0,\ 2$ & $3,\ 0$ \\ \cline{2-4}
\end{tabular}
\end{center}
Payoff vectors: player~1's rows over $(L,C,R)$ are $T = (2,1,4)$, $M = (3,1,2)$, $B = (1,0,3)$;
player~2's columns over $(T,M,B)$ are $L = (0,4,3)$, $C = (1,2,2)$, $R = (2,3,0)$.

\textbf{Round 1.} All pairs:
\begin{center}\small
\begin{tabular}{llll}
\toprule
Player & Pair & Componentwise comparison & Verdict \\
\midrule
1 & $T$ vs $M$ & $2<3$, $1=1$, $4>2$ & not ranked \\
1 & $T$ vs $B$ & $2>1$, $1>0$, $4>3$ & $B$ strictly dominated by $T$ \\
1 & $M$ vs $B$ & $3>1$, $1>0$, $2<3$ & not ranked \\
2 & $L$ vs $C$ & $0<1$, $4>2$, $3>2$ & not ranked \\
2 & $L$ vs $R$ & $0<2$, $4>3$, $3>0$ & not ranked \\
2 & $C$ vs $R$ & $1<2$, $2<3$, $2>0$ & not ranked \\
\bottomrule
\end{tabular}
\end{center}
$C$ is \emph{not} dominated yet: row $B$ is what keeps it alive ($2 > 0$ against $R$).
$S_1^1 = \{T, M\}$, $S_2^1 = \{L, C, R\}$.

\textbf{Round 2.} Player~2 on rows $(T, M)$: $L = (0,4)$, $C = (1,2)$, $R = (2,3)$.
$C$ vs $R$: $1 < 2$ and $2 < 3$, so $C$ is strictly dominated by $R$. $L$ vs $C$
($0<1$, $4>2$) and $L$ vs $R$ ($0<2$, $4>3$) are not ranked. Player~1 on all three columns:
$T$ vs $M$ unchanged, not ranked. $S_1^2 = \{T,M\}$, $S_2^2 = \{L, R\}$.

\textbf{Round 3.} Player~1 on $(L, R)$: $T = (2, 4)$, $M = (3, 2)$, not ranked. Player~2 on
$(T, M)$: $L = (0, 4)$, $R = (2, 3)$, not ranked. Stop:
\[
  S^\infty = \{T, M\} \times \{L, R\}.
\]

\textbf{Pure Nash equilibria: best-response marking of the whole matrix.}
\begin{center}
\begin{tabular}{r|c|c|c|}
\multicolumn{1}{r}{1 $\backslash$ 2} & \multicolumn{1}{c}{$L$} & \multicolumn{1}{c}{$C$} & \multicolumn{1}{c}{$R$}\\ \cline{2-4}
$T$ & $2,\ 0$ & $\br{1},\ 1$ & $\br{4},\ \br{2}$ \\ \cline{2-4}
$M$ & $\br{3},\ \br{4}$ & $\br{1},\ 2$ & $2,\ 3$ \\ \cline{2-4}
$B$ & $1,\ \br{3}$ & $0,\ 2$ & $3,\ 0$ \\ \cline{2-4}
\end{tabular}
\end{center}
Player~1: $BR_1(L) = \{M\}$ ($3 > 2 > 1$), $BR_1(C) = \{T, M\}$ (tie at 1), $BR_1(R) = \{T\}$
($4 > 3 > 2$). Player~2: $BR_2(T) = \{R\}$ ($2 > 1 > 0$), $BR_2(M) = \{L\}$ ($4 > 3 > 2$),
$BR_2(B) = \{L\}$ ($3 > 2 > 0$). Both payoffs are underlined only in $(T,R)$ and $(M,L)$. Cell by
cell inside $S^\infty$:
\begin{itemize}
  \item $(T, L)$: player~1 would switch to $M$ ($3 > 2$). Not NE.
  \item $(T, R)$: $T \in BR_1(R)$ and $R \in BR_2(T)$. NE, payoffs $(4,2)$.
  \item $(M, L)$: $M \in BR_1(L)$ and $L \in BR_2(M)$. NE, payoffs $(3,4)$.
  \item $(M, R)$: player~1 would switch to $T$ ($4 > 2$). Not NE.
\end{itemize}

\interp $C$ is a sensible reply only to beliefs that put weight on $B$; once player~2 knows that
player~1 is rational, $C$ goes (Lecture~2, slide~7). IESDS stops at four profiles because each
surviving strategy is a best reply to one of the opponent's surviving strategies. Nash
equilibrium adds the requirement that beliefs be correct and cuts the four to two. The two
equilibria are ranked differently by the players ($(T,R)$ favours player~1, $(M,L)$ player~2), so
the game has a coordination problem that neither IESDS nor Nash equilibrium resolves.

\verify \texttt{check\_A1.py} reproduces the rounds (round~1: $B$; round~2: $C$), checks that $C$
is not dominated in the full game, and finds by brute force over all nine cells that the pure Nash
equilibria are exactly $(T,R)$ and $(M,L)$. Gibbons gives no answer. Kumar's unofficial Gibbons
manual (Answer~1.2, not in \texttt{Livros/}) gives the same deletions and equilibria; it writes the
survivors as ``(T,M) and (L,R)'', meaning $\{T,M\}$ for player~1 and $\{L,R\}$ for player~2.

\mistakes
\begin{itemize}
  \item Deleting $C$ in round~1 because it looks poor. In the full game $C$ beats $R$ in row $B$
  and $L$ in row $T$; it is dominated only once $B$ is gone.
  \item Reading ``2.0'' as the number two and leaving player~2's payoff blank. It is the pair
  $(2,0)$; with any other reading the column comparisons change.
  \item Calling $(T,L)$ or $(M,R)$ an equilibrium because they survive IESDS. Survival is
  necessary for Nash equilibrium, not sufficient.
\end{itemize}

\answer{IESDS: delete $B$ (by $T$), then $C$ (by $R$); survivors $\{T, M\} \times \{L, R\}$.
Pure Nash equilibria: $(T, R)$ with payoffs $(4,2)$ and $(M, L)$ with payoffs $(3,4)$.}
\end{exsol}
